\subsection{Passes 3769--3786: the $GQ(4,2)$--Veldkamp--W33 closure}
\label{sec:bt3769-3786-gq-veldkamp}

The six-dimensional minus quadratic space over $\mathbb F_2$ supplies $27$ singular and $36$ nonsingular nonzero vectors.  Its singular points and totally singular lines form $GQ(2,4)$, with $27$ points and $45$ three-point lines; the dual $GQ(4,2)$ has $45$ points and $27$ five-point lines.  The dual point graph is
\[
 \operatorname{srg}(45,12,3,3),
 \qquad
 \operatorname{Spec}(A)=\{12^1,3^{20},(-3)^{24}\}.
\]
Writing $A_0=I$, $A_1=A$, and $A_2=J-I-A$, its Bose--Mesner multiplication is
\[
 A_1^2=12A_0+3A_1+3A_2,
 \quad
 A_1A_2=8A_1+9A_2,
 \quad
 A_2^2=32A_0+24A_1+22A_2.
\]
The primitive-idempotent ranks are $1,20,24$.  The Terwilliger algebra at a point has dimension $16$, center dimension $5$, and standard-module central block multiplicities $2,3,8,14,18$, hence semisimple type
\[
 M_3\oplus M_2\oplus \mathbb C^3.
\]

All $651$ projective lines of $PG(5,2)$ split by quadratic type as
\[
 45+216+270+120.
\]
The final $120$ lines consist of three nonsingular vectors and coincide exactly with the Norton--Fischer triples of the preceding spread algebra.  Thus that triple system is the type-IV line class of the Veldkamp space $V(GQ(2,4))$.

An exact-cover enumeration finds $200$ ovoids of $GQ(4,2)$, in automorphism orbits $40+160$.  On the $40$ plane ovoids, declare two objects adjacent when they share three quadrangle lines.  The resulting graph is
\[
 \operatorname{srg}(40,12,2,4)
\]
and is explicitly isomorphic to the independently reconstructed W33 point graph.  If $M$ is the $40\times45$ plane-ovoid incidence matrix, then
\[
 MM^{\mathsf T}=8I+2A_{W33}+J,
 \qquad
 \operatorname{Spec}(MM^{\mathsf T})=\{72^1,12^{24},0^{15}\}.
\]

The recurring count $135=45\cdot3$ supports two inequivalent $U_4(2)$ actions.  The Lagrangian-frame stabilizer has order census
\[
 1^1\,2^{43}\,3^{32}\,4^{84}\,6^{32},
\]
whereas an ordinary $GQ(4,2)$ flag stabilizer has
\[
 1^1\,2^{27}\,3^{32}\,4^{36}\,6^{96}.
\]
The corresponding permutation-character norms are $9$ and $8$, with cross-inner product $4$.

The rank-$24$ primitive projector is
\[
 P=(A-12I)(A-3I)=90E_{-3},
\]
with entries $48$ on the diagonal, $-12$ on adjacent pairs, and $3$ on nonadjacent pairs.  On $\operatorname{im}P$ define
\[
 x\star y=E_{-3}(x\circ y),
 \qquad
 a_i=Pe_i/21.
\]
Then the $45$ vectors $a_i$ are idempotent axes with Peirce spectrum
\[
 1^1,\qquad (1/7)^{14},\qquad (-1/3)^9.
\]
For collinear axes,
\[
 a_i\star a_j=-\frac13(a_i+a_j),
\]
and the five axes on every quadrangle line sum to zero.  For noncollinear $i,j$ with common neighbors $c_1,c_2,c_3$,
\[
 a_i\star a_j=rac17(a_i+a_j)+\frac5{42}(a_{c_1}+a_{c_2}+a_{c_3}).
\]
The algebra is positive Frobenius; its multiplication operators span all $M_{24}$ modulo $1000003$, so it is simple, and its derivation algebra is zero because $1/2$ is absent from every axis spectrum.  Its fusion law is not nontrivially $\mathbb Z_2$-graded, so this is not a Monster-$2A$ Majorana algebra.

For a representative $W(D_4)$ frame stabilizer, the invariant alternating forms modulo two have dimension seven and rank distribution
\[
 0^1,4^3,6^4,8^{12},10^{12},14^{40},16^{56}.
\]
No invariant form has rank $24$; consequently no $W(D_4)$-invariant even unimodular rank-$24$ child exists.  A deterministic trace-eight involution nevertheless permits an integral determinant-$-1$ basis change that pulls back $E_8^3$ to an even unimodular positive child whose exact surviving $U_4(2)$ stabilizer is $C_2$.  This child has roots and is not promoted as a Leech realization.

The finite-group embedding front remains fail-closed.  A future candidate must reproduce the complete $36/45/27/120/135/40$ fingerprint, independent order $51840$, class correspondence, and content-addressed incidence digests before any Monster embedding is asserted.
